\chapter{Desarrollo}
\label{ch:desarrollo}
% El desarrollo software, sin volcar código entero. Fuente: src/ + README.
% Frontera con metodología: aquí solo el CÓMO. El QUÉ/POR QUÉ está en \ref{ch:metodologia}.
% No repetir decisiones de diseño, sino describir su realización.
% Pendiente: el diagrama del pipeline de \ref{sec:pipeline}, cuando exista el estilo de figuras.

Este capítulo describe cómo se realizó en software el diseño del Capítulo~\ref{ch:metodologia}. No repite las decisiones de diseño, sino que expone su realización y las restricciones que esta impuso. La Sección~\ref{sec:pipeline} describe la estructura del programa que ejecuta un entrenamiento. La Sección~\ref{sec:des-metricas} detalla cómo se calculan las métricas y cuál es la optimización que hace viable medirlas en cada \emph{epoch} de los 960 entrenamientos. La Sección~\ref{sec:reproducibilidad} recoge las medidas de reproducibilidad y sus límites. Las Secciones~\ref{sec:ejecucion} y~\ref{sec:pilot} describen los dos lanzadores, el de la matriz completa y el del piloto de calibración. La Sección~\ref{sec:des-conclusiones} recapitula lo esencial.

\section{Arquitectura del pipeline}
\label{sec:pipeline}

Un entrenamiento es una configuración y un solo programa. La configuración se resuelve por capas, de menor a mayor prioridad: los valores por defecto, el fichero de la celda y las opciones de línea de órdenes, que son las que inyectan el \emph{learning rate} y la \emph{seed} de cada punto de la rejilla. A partir de ahí el programa fija las semillas, construye los cargadores de datos y el \emph{batch} de medición, construye la red y entrena hasta agotar el presupuesto de \emph{epochs}. Al terminar cada \emph{epoch} evalúa la validación, pone la red en modo de evaluación, calcula todas las métricas sobre el \emph{batch} de medición y registra una fila. Dos de esas elecciones son de realización y fijan qué se mide. La primera es medir en modo de evaluación. Para la red \emph{fully connected} y la convolucional no cambia nada, y en la ResNet-18, la única de las tres con \emph{batch normalization}, es la única forma de definir un gradiente por ejemplo, porque en modo de entrenamiento la salida de cada ejemplo depende de los demás ejemplos del lote; lo que se mide en ella es la red con sus estadísticas acumuladas, que en las primeras \emph{epochs} todavía no coinciden con las del lote. La segunda es la cadencia, una medición al final de cada \emph{epoch}, con lo que entre dos lecturas median entre 351 y 703 pasos del optimizador según el conjunto de datos. Esa cadencia es la causa directa del solape entre la ventana y el desenlace que acota la Sección~\ref{sec:res-solape}.

Una decisión de estructura sostiene el resto: los ejes de la matriz de experimentos viven en el mismo módulo que los ajustes de un entrenamiento individual. Los conjuntos de datos, las arquitecturas, los optimizadores, las rejillas de \emph{learning rate}, las \emph{seeds}, los presupuestos por conjunto y los ajustes congelados son constantes de ese módulo, y el resto del código las importa en lugar de repetirlas. El diseño congelado tiene así una definición única, y cambiar un eje es cambiar una línea en un sitio en vez de rastrear el mismo número por varios ficheros.

Cada entrenamiento deja cuatro ficheros en su propio directorio: la configuración ya resuelta, la trayectoria con una fila por \emph{epoch}, esa misma trayectoria recortada a las cinco ventanas, y un resumen con los indicadores de eficiencia, la única evaluación de test y los relojes. El resumen se escribe el último porque su existencia es lo que marca un entrenamiento como terminado, y de ella depende la reanudación que describe la Sección~\ref{sec:ejecucion}. Dos precisiones sobre lo que esos ficheros contienen. La columna de \emph{loss} de entrenamiento de la trayectoria es el \emph{loss} del último \emph{minibatch} de la \emph{epoch}, no la media de la \emph{epoch}; la media, que es lo que consume el estimador de velocidad de entrenamiento, se calcula aparte a partir de todos los pasos. Cuatro campos del resumen quedaron además obsoletos después de la matriz: las \emph{epochs} y los segundos hasta el umbral se calcularon con los umbrales de junio, y el mejor \emph{accuracy} y el mejor \emph{loss} de validación con un suavizado que en la primera y en la última \emph{epoch} tomaba la media de dos valores. La capa de análisis no los lee. Los recalcula desde la trayectoria con los umbrales definitivos de la Sección~\ref{sec:calibracion} y con el suavizado de la Sección~\ref{sec:variables}, que conserva el valor crudo en los bordes.

El módulo de registro se limita a la entrada y salida. Acumula filas y las escribe. Toda la lógica del dominio, el recorte de ventanas y el cálculo de los indicadores de eficiencia, queda en el programa de entrenamiento. La separación mantiene el registro trivial de leer y evita que una decisión de diseño acabe escondida en el código que guarda ficheros.

\section{Cálculo de las métricas}
\label{sec:des-metricas}

Las ocho métricas de gradiente comparten una interfaz única: reciben el modelo congelado y el \emph{batch} de medición, y devuelven un diccionario de escalares. Esa uniformidad es lo que permite recorrerlas en un bucle y añadir una novena sin tocar el bucle. El predictor de referencia queda deliberadamente fuera de ese registro, porque su firma es distinta: no recibe un modelo, sino la sucesión de \emph{loss} medios de entrenamiento por \emph{epoch}. Mantenerlos separados evita forzar una interfaz común entre dos cosas que no la comparten.

Cada métrica se parte en dos piezas. Un núcleo puro opera sobre tensores de gradiente ya extraídos y no sabe nada de modelos ni de datos, y una envoltura fina extrae los gradientes y llama a ese núcleo. La ventaja es de comprobación: el núcleo se puede alimentar con matrices de gradiente construidas a mano, cuyo resultado se conoce de antemano por cálculo directo, en lugar de contrastar la métrica contra sí misma sobre una red real. Las pruebas del repositorio se apoyan en esa separación y cubren cada núcleo con casos analíticos, más una prueba de humo sobre una red diminuta.

La optimización que hace viable el estudio es la del barrido de gradientes por ejemplo. Calcular por separado el gradiente de cada ejemplo del \emph{batch} de medición es, con diferencia, el coste dominante de una medición, y seis de las ocho métricas lo necesitan. En lugar de que cada una lo recalcule, el barrido se ejecuta una sola vez por medición y entrega todos los productos que esas seis consumen: los dos primeros momentos por columna, la matriz de productos escalares entre pares de gradientes junto con sus normas, y los cosenos de última capa. Cada métrica reduce entonces esos productos a sus escalares. Las dos restantes, la \emph{gradient disparity} y la varianza normalizada, trabajan con el gradiente medio de submuestras del mismo \emph{batch} de medición, cinco de 51 ejemplos y diez de 25. Ese gradiente medio es la suma de las filas de la submuestra, así que su norma y sus distancias son sumas de bloques de la matriz de productos escalares que el barrido ya calcula, y las dos métricas podrían leerse de él sin retropropagar nada más. Se calculan aparte, con sus propias pasadas, por fidelidad a la forma en que sus artículos las definen.

El segundo problema es de memoria. La matriz densa de gradientes por ejemplo tiene una fila por ejemplo del \emph{batch} de medición y una columna por parámetro, de modo que para la red \emph{fully connected} sobre Tiny-ImageNet ocupa 14,2 GB en precisión simple, más de lo que la GPU empleada deja libre una vez descontados el modelo y sus activaciones (\textbf{[pendiente: memoria exacta de la GPU, copiar de nvidia-smi]}; el equipo se describe en el Apéndice~\ref{sec:hardware}). En la GPU nunca se materializa entera: los gradientes se calculan por bloques de filas, en el acelerador reside un solo bloque cada vez, y de cada bloque se acumulan en el acto los momentos por columna y los cosenos de última capa. La matriz de productos escalares sí necesita todas las filas a la vez, y por eso cada bloque viaja a la memoria del anfitrión, donde la matriz completa se reúne en precisión simple, y vuelve al acelerador por bloques de columnas para formar los productos. El anfitrión tiene que disponer de esos 14 GB, y el tráfico de ida y vuelta, cuatro bytes por cada ejemplo y parámetro, es parte del coste real de la medición. El tamaño de bloque es un ajuste operativo y no científico: los estadísticos que salen del barrido son iguales para cualquier tamaño de bloque salvo redondeo, la prueba del repositorio los exige iguales hasta una parte por millón, y lo único que cambia es el pico de memoria.

El coste de una medición se descompone en dos cantidades que crecen de forma distinta: el número de ejemplos que hay que retropropagar, y la aritmética que se hace después sobre los gradientes ya calculados. Retropropagar un ejemplo es ejecutar sobre él la retropropagación, la operación que devuelve el gradiente de su \emph{loss} respecto a cada peso de la red, y es la misma que ejecuta cada paso de entrenamiento sobre un \emph{minibatch} entero. Las dos cantidades se cuentan aparte porque retropropagar un ejemplo no cuesta lo mismo en las tres arquitecturas. En la red \emph{fully connected} cuesta del orden del número de parámetros, mientras que en las convolucionales lo que fija el coste es el trabajo de las convoluciones, de modo que las dos cantidades no se pueden sumar en una sola cifra.

La Tabla~\ref{tab:coste-metricas} recoge ambas en los dos regímenes, cada métrica medida por separado y cada métrica añadida cuando el barrido compartido ya se ha ejecutado. En la columna de ejemplos retropropagados, un valor de $M$ indica que la métrica necesita el gradiente de cada ejemplo por separado. No son $M$ retropropagaciones sueltas: la herramienta de PyTorch que vectoriza el gradiente de una función sobre un lote, \texttt{vmap} aplicado a \texttt{grad} en \texttt{torch.func}, devuelve los gradientes de todo un bloque de filas en una sola llamada, con el coste de derivar cada ejemplo pero sin el bucle. Un número pequeño indica que a la métrica le basta el gradiente medio de cada submuestra, así que retropropaga una sola vez por submuestra, sobre todos sus ejemplos a la vez. Un cero indica que el barrido compartido ya hizo ese trabajo y la métrica se limita a reutilizar lo que aquel produjo; la primera fila de la tabla es ese barrido, con la aritmética que hace una sola vez por medición.

\begin{table}[htbp]
\centering
\footnotesize
\setlength{\tabcolsep}{3pt}
\caption{Coste asintótico de una medición, con los ejemplos retropropagados contados aparte de la aritmética posterior. $M$ es el tamaño del \emph{batch} de medición, $P$ el número de parámetros de la red, $W$ los pesos de su última capa y $T$ las \emph{epochs} transcurridas.}
\label{tab:coste-metricas}
\begin{tabular}{>{\raggedright\arraybackslash}p{3.5cm}
                >{\raggedright\arraybackslash}p{2.6cm}
                >{\raggedright\arraybackslash}p{2.0cm}
                >{\raggedright\arraybackslash}p{2.6cm}
                >{\raggedright\arraybackslash}p{2.2cm}}
\toprule
& \multicolumn{2}{l}{\textbf{Medida por separado}} & \multicolumn{2}{l}{\textbf{Con barrido compartido}} \\
\cmidrule(lr){2-3} \cmidrule(lr){4-5}
\textbf{Métrica} & \textbf{Ejemplos retropropagados} & \textbf{Aritmética} & \textbf{Ejemplos retropropagados} & \textbf{Aritmética} \\
\midrule
\multicolumn{5}{l}{\itshape Barrido compartido, una vez por medición} \\
\addlinespace[2pt]
Gradientes por ejemplo & \multicolumn{2}{l}{no se mide por separado} & $M$, una por ejemplo & $\Theta(M^2P + MP + MW)$ \\
\addlinespace[3pt]
\multicolumn{5}{l}{\itshape Familia de alineación o coherencia} \\
\addlinespace[2pt]
m-coherencia & $M$, una por ejemplo & $\Theta(MP)$ & 0 & $\Theta(P)$ \\
\addlinespace[2pt]
Stiffness & $M$, una por ejemplo & $\Theta(M^2P)$ & 0 & $\Theta(M^2)$ \\
\addlinespace[2pt]
Gradient confusion & $M$, una por ejemplo & $\Theta(M^2P)$ & 0 & $\Theta(M^2\log M)$ \\
\addlinespace[2pt]
Gradient disparity & 5, una por submuestra & $\Theta(P)$ & 5, una por submuestra & $\Theta(P)$ \\
\addlinespace[2pt]
GWA & $M$, una por ejemplo & $\Theta(MW)$ & 0 & $\Theta(M)$ \\
\addlinespace[3pt]
\multicolumn{5}{l}{\itshape Familia de variabilidad estocástica} \\
\addlinespace[2pt]
Varianza normalizada & 10, una por submuestra & $\Theta(P)$ & 10, una por submuestra & $\Theta(P)$ \\
\addlinespace[2pt]
Escala de ruido (GNS) & $M$, una por ejemplo & $\Theta(MP)$ & 0 & $\Theta(P)$ \\
\addlinespace[2pt]
GSNR & $M$, una por ejemplo & $\Theta(MP)$ & 0 & $\Theta(P\log P)$ \\
\addlinespace[3pt]
\multicolumn{5}{l}{\itshape Predictor de referencia, sin calcular el gradiente} \\
\addlinespace[2pt]
TSE & 0 & $\Theta(T)$ & 0 & $\Theta(T)$ \\
\bottomrule
\end{tabular}
\end{table}

El número de ejemplos retropropagados no separa a ninguna métrica de otra. Las ocho recorren el \emph{batch} de medición una vez, seis ejemplo a ejemplo y dos por submuestras, y entre esas dos formas la diferencia está en la memoria, porque el número de ejemplos derivados es prácticamente el mismo. Lo que sí las separa es la aritmética posterior, en dos niveles: las que acumulan momentos por columna crecen con el producto del tamaño del \emph{batch} por el número de parámetros, y las que necesitan todos los pares crecen con el cuadrado del tamaño del \emph{batch}. La segunda mitad de la tabla es la que justifica el barrido compartido. Seis de las ocho métricas no necesitan ni una sola pasada más, y solo las dos que trabajan con submuestras siguen pagando las suyas, con lo que medir el registro completo cuesta unas tres pasadas sobre el \emph{batch} de medición en lugar de ocho. Lo que domina el tiempo real no son esas pasadas, sino la aritmética posterior y el tráfico de la matriz con el anfitrión.

La GWA solo lee los pesos de la última capa, pero la implementación deriva la red entera para cada ejemplo y descarta después el resto del gradiente, así que su ventaja está en la aritmética y no en el número de ejemplos retropropagados; esa última capa, además, no siempre es pequeña, porque en la red convolucional sobre Tiny-ImageNet reúne 102.400 de los 116.936 pesos.

Los momentos se acumulan en doble precisión, porque las formas que restan cuadrados sufren cancelación en precisión simple. Una métrica que falle, por ejemplo al agotar la memoria en la matriz de pares, se registra y se omite: sus columnas quedan ausentes de esa fila y el entrenamiento continúa. Un fallo aislado cuesta unas columnas, no un entrenamiento de horas. En la matriz no ocurrió ninguna vez. Las 33.600 filas de trayectoria de los 960 entrenamientos llevan las 23 columnas de métrica de gradiente y las cuatro del predictor de referencia, y en las filas de entrenamientos vivos en esa \emph{epoch}, con \emph{loss} finito y red sin colapsar, ninguna de las 27 tiene un valor ausente; los ausentes que hay están todos en filas de entrenamientos divergidos o colapsados.

Lo que esa aritmética cuesta de verdad se lee en los relojes de los 960 resúmenes, que separan el tiempo de medición del resto. Sumado sobre los 40 entrenamientos de cada celda, el tiempo de medición va del 0,7\,\% del tiempo de entrenamiento en la red convolucional sobre Tiny-ImageNet al 105\,\% en la red \emph{fully connected} sobre CIFAR-100 con SGD, donde la celda cuesta 2,05 veces lo que costaría sin instrumentar; por arquitectura, la convolucional queda por debajo del 3\,\%, la ResNet-18 entre el 8 y el 65\,\% y la red plana entre el 38 y el 105\,\%. La red más barata de entrenar es la que paga la instrumentación más cara, porque el coste de la medición crece con el número de parámetros, que en la red plana es el mayor de las tres, mientras que su paso de entrenamiento es el más ligero. En total, de las 121,7 horas de la matriz, 97,6 fueron entrenamiento y 24,1 medición, una cuarta parte del tiempo de entrenamiento.

\section{Reproducibilidad}
\label{sec:reproducibilidad}

Que dos entrenamientos de una celda solo difieran en el \emph{learning rate} y en la \emph{seed} lo garantiza la siembra, no el modo determinista. Al arrancar, cada entrenamiento fija la semilla de los tres generadores en juego, el de Python, el de NumPy y el de PyTorch, y con ella quedan fijados la inicialización, el orden de los \emph{minibatches} y el \emph{batch} de medición. Lo que añade el modo determinista de las bibliotecas de convolución, que renuncia a la selección automática de algoritmos, es que un mismo entrenamiento se repita bit a bit en la misma máquina. Ese alcance es acotado. Se fija el modo determinista de cuDNN y se desactiva su selección de algoritmos, pero no se activa el modo determinista general de PyTorch, que rechazaría las operaciones sin variante determinista, y la semilla de \emph{hash} de Python que el programa escribe en su propio entorno no tiene efecto sobre un intérprete ya en marcha. La repetición bit a bit está comprobada sobre los datos de este trabajo. Veintidós de los 24 entrenamientos del piloto tienen su gemelo en la matriz, con el mismo conjunto de datos, arquitectura, optimizador, \emph{learning rate} y \emph{seed}, ejecutado dos meses después con otra versión del código, y en los 22 las curvas de \emph{loss} de entrenamiento, de \emph{loss} de validación y de \emph{accuracy} de validación coinciden bit a bit en las 20 o 40 \emph{epochs} comunes, y las 23 columnas de métrica coinciden con una diferencia relativa por debajo de $10^{-12}$. Los dos gemelos restantes, la red convolucional sobre MNIST con los dos optimizadores, difieren desde la primera \emph{epoch} en menos de dos milésimas de \emph{accuracy}, y la causa de esa diferencia no está establecida.

Hay dos semillas distintas. La de partición es fija e idéntica en los 960 entrenamientos, de modo que todos ven el mismo reparto entre entrenamiento y validación. La del entrenamiento gobierna la inicialización de los pesos, el orden de los \emph{minibatches} y qué ejemplos forman el \emph{batch} de medición, que son las fuentes de variación que el estudio sí quiere observar.

El \emph{batch} de medición se construye una sola vez, al empezar, y queda como un par de tensores fijos que se reutilizan en todas las mediciones. Así la métrica de la primera \emph{epoch} y la de la última se calculan sobre exactamente los mismos ejemplos, y la trayectoria refleja el cambio de la red y no el de la muestra. Esos ejemplos siguen en el conjunto de entrenamiento y la red sigue aprendiendo de ellos, como en los artículos de origen de las métricas, así que en las últimas \emph{epochs} las métricas describen el gradiente residual sobre ejemplos ya memorizados; una sonda tomada de la validación mediría otra cantidad.

El determinismo tiene un límite. Las constantes de normalización de cada conjunto de datos se recalcularon desde cero para comprobarlas: coinciden hasta $5 \times 10^{-5}$ en MNIST y en los dos CIFAR, pero solo hasta unos $6 \times 10^{-4}$ en Tiny-ImageNet, con la media exacta y la desviación típica algo menor en los tres canales. La diferencia es compatible con versiones distintas de la biblioteca que decodifica las imágenes JPEG, formato que Tiny-ImageNet emplea y los otros tres conjuntos no. Es una discrepancia pequeña y sin efecto sobre las comparaciones internas del estudio, porque los 960 entrenamientos usan las mismas constantes, pero implica que reproducir estas cifras en otra máquina puede exigir fijar también esa versión. Otra máquina cambia además el redondeo de todas las operaciones y, en Tiny-ImageNet, los propios píxeles, así que lo que se reproduce fuera de esta máquina es la estadística de las 24 celdas y no la trayectoria de cada entrenamiento. El entorno completo, con versiones y hardware, se detalla en el Apéndice~\ref{ap:configuracion}.

\section{Ejecución y reanudación}
\label{sec:ejecucion}

La matriz se ejecuta con un lanzador que hace tres cosas. Genera las 24 configuraciones de celda que falten, respetando las que ya existan para que una corrección a mano sobreviva. Informa del estado, es decir, cuántos puntos de la rejilla están hechos y cuántos pendientes. Ejecuta los pendientes, cada uno como un proceso independiente.

La reanudación no necesita ningún registro externo. El nombre de un entrenamiento queda determinado por sus cinco coordenadas, arquitectura, conjunto de datos, optimizador, \emph{learning rate} y \emph{seed}, así que el lanzador puede construir el nombre de cualquier punto de la rejilla y comprobar si su resumen existe. Si existe, el punto está hecho. Volver a lanzar la misma orden después de una interrupción continúa por donde iba, y lanzarla sobre una matriz completa no ejecuta nada. La propiedad depende por entero de que el resumen se escriba el último, como señala la Sección~\ref{sec:pipeline}. El registro acumula las filas en memoria y escribe la trayectoria después de la última \emph{epoch}, así que un entrenamiento interrumpido a mitad deja solo su configuración, ni trayectoria ni resumen, y se repite entero; no hay puntos de control dentro de un entrenamiento.

El lanzador admite además restringirse a una parte de la rejilla, por conjunto de datos, arquitectura u optimizador. El estudio se ejecutó en una sola GPU y sin acceso a un clúster, de modo que ese troceado sirve para avanzar la matriz por tramos en sesiones distintas, no para repartir el trabajo entre máquinas. Cada entrenamiento corre como un proceso aparte con el código que hay en disco al lanzarlo, así que lo que garantiza que los 960 usaran el mismo código es el historial del repositorio: entre el primer entrenamiento versionado, del 8 de agosto de 2026, y el último, del 22 de agosto, ni el programa de entrenamiento ni el paquete de métricas recibieron ningún cambio.

\section{Piloto de calibración}
\label{sec:pilot}

El piloto de calibración tiene su propio lanzador, que reutiliza el de la matriz. Ejecuta un entrenamiento por celda, con el \emph{learning rate} en el centro de la rejilla de su optimizador, \emph{seed} 0 y el doble del presupuesto candidato, y ofrece después una lectura tabulada de las curvas obtenidas. El protocolo y los valores que de ahí salieron están en la Sección~\ref{sec:calibracion}. La regla codificada que producía la recomendación de presupuesto era esta: la meseta del \emph{loss} de validación es la primera \emph{epoch} cuyo \emph{loss} queda a menos del 2\,\% del mínimo del entrenamiento, y el presupuesto recomendado es la meseta más alta del conjunto de datos con un margen del 20\,\%, redondeada hacia arriba al múltiplo de 20 siguiente. Lo que se adoptó frente a lo que la regla recomendó lo cuenta la Sección~\ref{sec:calibracion}.

Dos cosas separan el piloto de la matriz. La primera es el código. El piloto corrió a mediados de junio de 2026, y desde entonces cambiaron el cargador de Tiny-ImageNet, cuyo etiquetado del conjunto de test tenía un error corregido el 18 de junio, y la asignación de memoria del acelerador, mientras que el barrido compartido de la Sección~\ref{sec:des-metricas} entró en esos mismos días. Los resúmenes del piloto sobre Tiny-ImageNet declaran en un campo propio que su \emph{accuracy} de test y su \emph{gap} se calcularon con el error, y cada uno guarda dentro una repetición corregida con el presupuesto definitivo de 40 \emph{epochs}; la calibración solo usó las curvas de validación y los relojes, que no dependen del etiquetado del test. La segunda es la lectura. El informe del piloto lee hoy el presupuesto y el umbral del módulo de configuración, que ya contiene los valores definitivos, así que la tabla que imprime no es la que se leyó en junio, y para reproducir aquella lectura hay que volver al estado del código de esa fecha.

El piloto escribe en un directorio distinto del de la matriz, y esa separación evita un error silencioso. El lanzador de la matriz da un punto por terminado cuando encuentra su resumen, así que un entrenamiento del piloto escrito en el directorio de la matriz se contaría más tarde como un punto hecho, cuando en realidad se entrenó con otro presupuesto.

\section{Conclusiones del capítulo}
\label{sec:des-conclusiones}

La realización software del diseño se apoya en cuatro decisiones. La primera es tener una definición única de la matriz, importada por todo lo demás, de modo que el diseño congelado no pueda divergir entre módulos. La segunda es la interfaz uniforme de las métricas, que separa un núcleo comprobable con casos analíticos de la extracción de gradientes. La tercera es el barrido compartido de gradientes por ejemplo, que amortiza el coste dominante de la medición entre las seis métricas que lo necesitan y que, junto al procesado por bloques de filas, es lo que permite medir en cada \emph{epoch} arquitecturas cuya matriz de gradientes por ejemplo no cabe en memoria. La cuarta es la reanudación basada en un único fichero marcador, que hace idempotente el lanzamiento de la matriz sin ningún estado externo.

Las cuatro cambian qué resulta posible medir con una sola GPU y en un tiempo acotado, que es la restricción real de este trabajo, y no lo que el Capítulo~\ref{ch:metodologia} fija como objeto de medida. Dos elecciones de realización sí lo tocan y quedan declaradas en la Sección~\ref{sec:pipeline}: medir con la red en modo de evaluación y medir una vez por \emph{epoch}.
